% BioDCASE 2026 Workshop Paper
% NOTE: Replace IEEEtran.cls with the official DCASE 2026 template when available.
%       Get IEEEtran.cls from https://www.ieee.org/conferences/publishing/templates.html
%       or run: texmf-update after installing texlive-publishers (Debian/Ubuntu).
%
% Compile: pdflatex main.tex && bibtex main && pdflatex main.tex && pdflatex main.tex

\documentclass[conference,a4paper]{IEEEtran}

\usepackage[utf8]{inputenc}
\usepackage[T1]{fontenc}
\usepackage{amsmath,amssymb}
\usepackage{graphicx}
\usepackage{booktabs}
\usepackage{multirow}
\usepackage{xcolor}
\usepackage{hyperref}
\usepackage{microtype}
\usepackage{algorithm}
\usepackage{algorithmic}
\usepackage{cite}
\usepackage{ifthen}

\hypersetup{colorlinks=false, pdfborder={0 0 0}}

% ------------------------------------------------------------------
% Convenience macros
% ------------------------------------------------------------------
\newcommand{\BAu}{BA_{\text{unseen}}}
\newcommand{\BAs}{BA_{\text{seen}}}
\newcommand{\DSG}{\mathrm{DSG}}
\newcommand{\best}[1]{\textbf{#1}}
\definecolor{pending}{rgb}{0.55,0.55,0.55}

% ------------------------------------------------------------------
\title{Domain-Balanced Adversarial Training with\\
Physics-Grounded Augmentation for\\
Cross-Domain Mosquito Species Classification}

\author{%
  \IEEEauthorblockN{Sulagna Saha}
  \IEEEauthorblockA{Mila -- Quebec AI Institute\\
    Université de Montréal\\
    \texttt{saha23s@mtholyoke.edu}}
  % TODO: add co-authors
}

\begin{document}

\maketitle

\begin{abstract}
Mosquito classifiers trained in one recording environment fail badly
when deployed in another---the domain-shift problem.
The BioDCASE 2026 challenge~\cite{biodcase2026} makes this concrete:
we find that one recording environment (D5) contains 99.4\% of all
training data, so standard domain-generalisation methods never
actually encounter the minority environments they need to generalise
to, and gain less than 0.5 percentage points on the true
cross-environment test.
The fix is simple but non-obvious: \textit{balance the training
domains first}.
Resampling to equalise domain frequency per epoch alone gains
+15.6~pp in unseen-domain balanced accuracy ($\BAu$).
With balanced training in place, domain-adversarial learning and
contrastive alignment across recording conditions push the system to
$\BAu{=}0.378$ (+21.3~pp) and a domain-generalisation gap of
$\DSG{=}0.202$ over the unbalanced baseline.
We also introduce \textit{Frequency-Band-Selective Mix~(FBS-Mix)}: an
augmentation that mixes the recording-condition-sensitive
low-frequency band of the spectrogram while leaving the mosquito
wingbeat band untouched, guided by a per-frequency variance analysis.
FBS-Mix requires no model changes and reduces $\DSG$ to 0.229 as a
lightweight alternative.
\end{abstract}

% ------------------------------------------------------------------
% =================================================================
\section{Introduction}
\label{sec:intro}
% =================================================================

Mosquitoes transmit diseases that kill hundreds of thousands of people
every year.
Knowing which species are present in a region---and in what numbers---is
the foundation of any targeted control effort, but traditional
identification requires physical capture and microscopy.
Miniaturised microphones can passively record the characteristic
wingbeat sound of flying mosquitoes, and machine learning can identify
the species from that sound alone~\cite{mukundarajan2017,kiskin2020humming}.
The practical barrier is \emph{domain shift}: a classifier trained on
recordings from one microphone or environment performs poorly when
tested on a different one, even for the same species.

The BioDCASE 2026 Cross-Domain Mosquito Species Classification
challenge~\cite{biodcase2026} formalises this problem.
Nine species are recorded across five conditions (D1--D5), and the
key metric is \emph{species balanced accuracy on unseen conditions}
($\BAu$)---how well the model recognises each species in a recording
environment it has never trained on.
A companion metric, the \emph{domain generalisation gap}
($\DSG{=}|\BAs - \BAu|$, lower is better), quantifies how much worse
the model gets on new conditions compared to known ones.

We diagnose a critical problem with the dataset that has not been
widely acknowledged.
One recording condition---D5---accounts for 212,339 of the 213,647
training clips: \textbf{99.4\%}.
The other four conditions together provide only 1,308 clips.
This extreme imbalance means that during training, a model almost
never sees examples from D1--D4.
Standard domain-generalisation methods like DANN~\cite{dann2016} and
MixStyle~\cite{mixstyle2021} are built on the assumption that training
data covers all domains; here, they do not.
Applied naively, they improve $\BAu$ by only 0.4 and 0.2 percentage
points respectively---essentially nothing.
The official challenge metric is similarly affected: roughly 85\% of
the test set looks like D5 by recording characteristics, so the
reported $\BAu{=}0.175$ mostly reflects D5 generalisation, not
genuine cross-condition transfer.
We adopt a Leave-One-Domain-Out~(LODO) evaluation instead, holding
out each of D1--D4 in turn, to measure true cross-condition
performance (Section~\ref{sec:data}).

Our contributions:
\begin{itemize}
  \item \textbf{A diagnostic finding:} We show that fixing the domain
    imbalance is a \emph{prerequisite} for domain-generalisation
    methods to work at all.
    Domain-balanced resampling alone gains +15.6~pp $\BAu$---the
    single largest improvement in our ablation.
  \item \textbf{FBS-Mix:} A data-motivated augmentation that mixes
    only the recording-noise frequency band
    (mel bins 0--8, $\approx$36--325~Hz) and leaves the wingbeat band
    (bins 9--35, $\approx$361--1{,}360~Hz) untouched, guided by a
    per-frequency variance analysis.
  \item \textbf{A comprehensive ablation:} We test ten method
    combinations under LODO, including actionable negative results:
    test-time batch-norm adaptation (TTBN) is harmful here, unbalanced
    DANN is near-useless, and a transformer backbone (AST) fails to
    benefit from adversarial training because its gradient reversal
    layer never suppresses domain information.
\end{itemize}

Our best system---balanced resampling, domain-adversarial
training~(DANN), and domain-invariant contrastive
learning~(DiCL)~\cite{drbiol2025} with a projection head---achieves
$\BAu{=}0.378$ (+21.3~pp over the unbalanced MTRCNN~\cite{mtrcnn2025}
baseline) and $\DSG{=}0.202$.
FBS-Mix gives a lighter-weight alternative: $\DSG{=}0.229$ with no
model changes.

% =================================================================
\section{Dataset \& Problem Analysis}
\label{sec:data}
% =================================================================

\subsection{Dataset}
\label{sec:data:dataset}

The CD-MSC corpus contains 271,380 audio clips of nine mosquito
species recorded across five conditions (D1--D5).
The labels D1--D5 are opaque---the challenge release does not
include device or location metadata---so we treat them as five distinct
but uncharacterised recording environments~\cite{biodcase2026}.
All clips are resampled to 8~kHz (the challenge-specified rate) and
converted to 64-bin log-mel spectrograms using the official baseline
frontend: 512-sample Hann window (25~ms), 80-sample hop (10~ms),
mel filterbank from 0 to 4{,}000~Hz.
Because mel frequency spacing is non-linear, the bin-to-frequency
mapping ranges from $\sim$36~Hz per bin at the low end to
$\sim$140~Hz per bin at the high end.
Features are normalised by subtracting the per-bin mean and dividing
by the per-bin standard deviation, both computed from the training
split only.
Clip lengths are mostly 0.63~s (mode = 63 frames); the test set
contains some multi-second recordings.
The official training/validation/test split is 213,647 / 30,516 /
27,217 clips.

\subsection{The Core Problem: One Domain Dominates Everything}
\label{sec:data:imbalance}

Table~\ref{tab:domain_counts} and Fig.~\ref{fig:data_analysis}~(left)
show the training-set breakdown.
D5 holds \textbf{212,339 of the 213,647 training clips---99.4\%}.
Conditions D1--D4 together contribute only 1,308 clips,
split across 80 to 634 clips each.

\begin{table}[t]
  \centering
  \caption{Training-set domain distribution. D5 holds 99.4\% of
           training data; D1--D4 together account for only 1,308 clips.}
  \label{tab:domain_counts}
  \setlength{\tabcolsep}{5pt}
  \begin{tabular}{lrrrr}
    \toprule
    Condition & Train & Val & Test & Train \% \\
    \midrule
    D1 & 634   & 96     & 3{,}335  & 0.30 \\
    D2 & 230   & 30     & 524      & 0.11 \\
    D3 & 364   & 53     & 262      & 0.17 \\
    D4 & 80    & 3      & 117      & 0.04 \\
    D5 & 212{,}339 & 30{,}334 & 22{,}979 & \textbf{99.38} \\
    \midrule
    Total & 213{,}647 & 30{,}516 & 27{,}217 & 100 \\
    \bottomrule
  \end{tabular}
\end{table}

\paragraph{Why this breaks domain-generalisation methods.}
Domain-adversarial training (DANN) works by training a discriminator
to classify which domain a sample came from, then reversing its
gradients to push the model toward domain-agnostic features.
But when 99.4\% of batches come from D5, the discriminator's easiest
strategy is to learn one thing: ``is this D5 or not?''
Reversing that gradient removes only the D5-vs-not-D5 signal, leaving
all D1--D4-specific features in the embedding intact.
When training data is balanced across domains, the discriminator
must separate all five conditions, and the reversed gradient removes
all of them.
This mechanism explains our central empirical finding: DANN without
balanced training gains only +0.4~pp $\BAu$; DANN \emph{with}
balanced training gains +20.3~pp.

\paragraph{Why D5 is excluded from our evaluation.}
Holding D5 out in a Leave-One-Domain-Out experiment would mean
training on only 1,308 clips---roughly the size of a small pilot
study---which is a data-starvation problem, not a domain-shift problem.
We exclude D5 from all LODO means; all reported numbers are D1--D4
mean.

\subsection{Why the Official Evaluation Metric Is Misleading}
\label{sec:data:split}

We compute the domain centroid of each condition in the training set
(the per-domain average log-mel feature vector) and assign each test
clip to its nearest centroid by Euclidean distance.
Approximately 85\% of test clips are closest to the D5 centroid.
The official 10-seed challenge baseline reports
$\BAu{=}0.175$~\cite{biodcase2026}---but 85\% of those ``unseen''
test clips effectively look like D5, making this a measure of
\emph{D5 generalisation}, not cross-condition transfer.
Our LODO baseline (seed~42, D1--D4 mean) gives
$\BAu{=}0.165$, $\BAs{=}0.578$, $\DSG{=}0.412$.
\textbf{We use LODO D1--D4 mean $\BAu$ as our primary metric
throughout.}

\subsection{Which Frequency Bands Carry Species vs.\ Recording-Condition Information?}
\label{sec:data:freq_analysis}

To decide \emph{where} in the spectrogram to apply data augmentation,
we ask: which mel bins vary more \emph{within} one recording condition
(species-specific variation) than \emph{across} conditions
(condition-specific variation)?
We compute this ratio---within-domain variance divided by
cross-domain variance---for each of the 64 mel bins, using the
held-out test split.\footnote{This analysis uses test-split data solely
to motivate the design choices $K{=}9$ and $K'{=}36$.
Both values were fixed before any FBS-Mix experiments; no further
tuning was performed.}

A ratio below~1 means the bin varies more across recording conditions
than within them: it is picking up \emph{recording noise}, not
species signal.
A ratio above~1 means it varies more within a condition: it is
picking up something that differs between species, not between
recording devices.

Fig.~\ref{fig:data_analysis}~(right) shows three clear zones:

\begin{itemize}
  \item \textbf{Bins 0--8} ($\approx$36--325~Hz), ratio 0.30--0.93:
    condition-dominated.
    These bins capture low-frequency recording noise and microphone
    response---information that changes with the device, not the mosquito.
  \item \textbf{Bins 9--35} ($\approx$361--1{,}360~Hz), ratio 1.50--4.12:
    species-dominated.
    This is the mosquito wingbeat core: the ratio peaks at bin~17
    ($\approx$649~Hz), near the wingbeat fundamental of the most
    common species~\cite{mukundarajan2017}.
  \item \textbf{Bins 36--63} ($\approx$1{,}412--3{,}854~Hz),
    ratio 0.73--1.47: mixed---upper harmonics and background noise
    are tangled together.
\end{itemize}

\textbf{Why mosquito wingbeat is device-invariant.}
Mosquito wings beat at a species-specific rate (300--800~Hz
fundamental, with harmonics extending to $\sim$1{,}400~Hz)
determined by wing morphology~\cite{mukundarajan2017}.
This physical oscillation is independent of the recording device, so
the wingbeat frequency content stays consistent across conditions even
when everything else changes.

This analysis tells us precisely what MixStyle gets wrong.
MixStyle~\cite{mixstyle2021} mixes the statistics of \emph{all}
frequency channels inside the CNN, which randomises the wingbeat
core (bins 9--35) at the same time as the condition-noise band
(bins 0--8).
Corrupting species information cancels any benefit from diversifying
condition noise---which is why balanced+MixStyle and balanced-only
perform identically ($\approx$0.32~$\BAu$).
FBS-Mix~(Section~\ref{sec:method:fbsmix}) mixes only the
condition-dominated band and leaves the wingbeat core alone.

% Figures for Sec 2 inserted here so they float near first reference
% =================================================================
% Figure definitions — included from main.tex between sections
% =================================================================

% ── Fig 1 ────────────────────────────────────────────────────────────────────
\begin{figure}[t]
  \centering
  \includegraphics[width=\columnwidth]{figures/fig1_data_analysis}
  \caption{%
    \textbf{Dataset imbalance and per-frequency variance analysis.}
    \textit{Left:} Training-set domain distribution (log scale).
    Domain D5 holds 212,339 clips~(99.4\%); D1--D4 together have only
    1,308.
    When D5 is the only available domain signal, adversarial training
    learns to detect ``D5 vs.\ not-D5'', removing only this one
    discriminative direction from the embedding space.
    \textit{Right:} Per-mel-bin ratio of within-domain to cross-domain
    variance on the test split.
    Bins~9--35 ($\approx$361--1{,}360~Hz, the wingbeat fundamental and
    lower harmonics) are consistently species-dominated;
    bins~0--8 ($\approx$36--325~Hz) are domain-dominated.
    FBS-Mix mixes only bins~0--8.%
  }
  \label{fig:data_analysis}
\end{figure}

% ── Fig 2 ────────────────────────────────────────────────────────────────────
\begin{figure}[t]
  \centering
  \includegraphics[width=\columnwidth]{figures/fig2_system_diagram}
  \caption{%
    \textbf{Proposed system.}
    FBS-Mix randomises the domain-dominated low-frequency band
    (bins~0--8, orange) before feature extraction, leaving the
    species-discriminative wingbeat core (bins~9--35, blue) unchanged.
    MTRCNN extracts a 32-dim shared embedding~$z$.
    DANN enforces domain invariance via gradient reversal (GRL).
    DiCL pulls same-species, cross-domain embeddings together through a
    projection MLP~($32{\to}128{\to}128$).%
  }
  \label{fig:system}
\end{figure}

% ── Fig 3 ────────────────────────────────────────────────────────────────────
\begin{figure}[t]
  \centering
  % TODO: replace with fig3_tsne.pdf once generated on compute node.
  % Run: python scripts/paper/fig3_tsne.py --baseline <path> --best <path>
  \IfFileExists{figures/fig3_tsne.pdf}{%
    \includegraphics[width=\columnwidth]{figures/fig3_tsne}
  }{%
    \fbox{\parbox{\columnwidth}{\centering\vspace{1.8em}
      \textit{[Placeholder: t-SNE figure pending GPU computation.}\\
      \textit{Run \texttt{scripts/paper/fig3\_tsne.py} on a compute node.]}
      \vspace{1.8em}
    }}%
  }
  \caption{%
    \textbf{t-SNE of 32-dim embeddings (LODO D1 test set, coloured by domain)
    --- pending GPU computation; figure will be populated for camera-ready.}
    \textit{Expected left (baseline):} D1--D4 embeddings collapse to an
    isolated corner, completely separated from the D5 cloud (confirmed on
    standard-split in prior analysis).
    \textit{Expected right (proposed system):} Domain clusters intermix
    substantially after adversarial and contrastive domain-invariance training.%
  }
  \label{fig:tsne}
\end{figure}

% =================================================================
\section{Proposed System}
\label{sec:method}
% =================================================================

Fig.~\ref{fig:system} shows the full pipeline.
We use the MTRCNN backbone~\cite{mtrcnn2025} throughout: three parallel
convolutional branches (temporal kernel sizes 3, 5, 7 capture different
wingbeat timescales), each with three Conv-BN-ReLU-Pool stages,
followed by masked mean$+$max pooling that ignores padding frames.
The three branch outputs are concatenated into a 32-dimensional
shared embedding~$z$ ($\approx$221K parameters total; the 32-dim size
follows the original MTRCNN paper).
Two linear heads predict species (9 classes) and recording domain
(5 classes) from~$z$.

\subsection{Domain-Balanced Sampling}
\label{sec:method:balance}

\textbf{The problem it solves:} without intervention, D5 appears in
99.4\% of every training batch, so the model learns almost exclusively
from one recording condition.
The other four conditions are effectively invisible.

\textbf{The fix:} we assign each training clip a sampling weight
$w_i = 1 / \text{count}(d_i)$, where $d_i$ is its domain label.
This equalises how often each domain is seen per epoch, upsampling
D1--D4 clips by 200--2{,}600$\times$.
As discussed in Section~\ref{sec:data:imbalance}, every subsequent
method in this paper depends on this step: without balanced training,
the model has no D1--D4 signal to generalise from.

\subsection{Domain-Adversarial Training (DANN)}
\label{sec:method:dann}

\textbf{The problem it solves:} even with balanced training, the
embedding~$z$ may still encode recording-condition information
alongside species information, letting the model ``cheat'' by
memorising condition-specific cues.

\textbf{How it works:} a Gradient Reversal Layer~(GRL)~\cite{dann2016}
sits between~$z$ and the domain head.
During the forward pass it is an identity; during backpropagation it
flips the sign and scales the gradient by~$\lambda(p)$.
This makes the embedding \emph{worse} at predicting the domain---the
model is forced to discard recording-condition information to keep the
domain head from working.
The schedule from~\cite{dann2016} gradually increases the adversarial
pressure:
\begin{equation}
  \lambda(p) = \lambda_{\max}\!\cdot\!\left(\frac{2}{1+e^{-10p}}-1\right),
  \quad p = \frac{e}{E},
  \label{eq:dann_schedule}
\end{equation}
where $e$ is the current epoch and $E{=}100$ (max epochs).
The constant~10 controls how quickly the adversary ramps up~\cite{dann2016};
at the halfway point ($p{=}0.5$) the model is already at 96\% of its
full adversarial pressure.
We set $\lambda_{\max}{=}1.0$.

\subsection{Domain-Invariant Contrastive Learning (DiCL)}
\label{sec:method:dicl}

\textbf{The problem it solves:} DANN pushes domain information
\emph{out} of the embedding, but it does not explicitly pull
same-species recordings from different conditions
\emph{together}.

\textbf{How it works:} DiCL~\cite{drbiol2025} treats any two clips of
the same species from \emph{different} recording conditions as a
positive pair.
For each such anchor~$i$, the contrastive loss pulls its embedding
toward all its cross-condition same-species partners, while pushing
it away from all other clips in the batch.

Concretely, let $P_i = \{j : y_j = y_i,\, d_j \neq d_i\}$ be the
cross-condition positives for anchor~$i$, and let
$\mathcal{A} = \{i : |P_i| > 0\}$.
Embeddings are first passed through a two-layer projection head
($h_i = \mathrm{proj}(z_i) \in \mathbb{R}^{128}$; architecture:
$32{\to}128{\xrightarrow{\mathrm{ReLU}}}128$, no final activation).
The projection head means the contrastive geometry is learned in a
separate space from the classifier, which prevents the loss from
distorting the species-discriminative directions in~$z$~\cite{supcon2020}.
DiCL then minimises:
\begin{equation}
  \mathcal{L}_{\mathrm{DiCL}} =
  {-}\frac{1}{|\mathcal{A}|}\!\sum_{i \in \mathcal{A}}
  \frac{1}{|P_i|}\!\sum_{j \in P_i}
  \log\frac{e^{\,\tilde{h}_i \cdot \tilde{h}_j / \tau}}
           {\sum_{k \neq i} e^{\,\tilde{h}_i \cdot \tilde{h}_k / \tau}},
  \label{eq:dicl}
\end{equation}
where $\tilde{h} = h / \|h\|_2$.
We use temperature $\tau{=}0.2$ (rather than 0.07 in~\cite{drbiol2025}):
with only 80--634 minority-domain clips per epoch there are very few
positive pairs per batch, and a lower temperature concentrates the
gradient on only the hardest negatives, making learning unstable.
The combined training loss is:
\begin{equation}
  \mathcal{L} = \mathcal{L}_{\mathrm{sp}} - \lambda(p)\,\mathcal{L}_{\mathrm{dom}}
              + w_c\,\mathcal{L}_{\mathrm{DiCL}},
  \label{eq:total_loss}
\end{equation}
with $w_c{=}0.1$~\cite{drbiol2025} and cross-entropy losses
$\mathcal{L}_{\mathrm{sp}}$, $\mathcal{L}_{\mathrm{dom}}$ for species
and domain.

\subsection{Frequency-Band-Selective Mix (FBS-Mix)}
\label{sec:method:fbsmix}

\textbf{The problem it solves:}
MixStyle~\cite{mixstyle2021} diversifies domain characteristics by
randomly blending the mean and variance of CNN feature maps between
training samples.
It works well on vision tasks, but here it backfires: it mixes the
frequency bands that carry \emph{species} information (the wingbeat
tone) at the same time as the recording-noise bands it is trying to
diversify.
The two effects cancel, and balanced+MixStyle performs no better than
balanced alone ($\approx$0.32~$\BAu$).

\textbf{The insight:}
the per-frequency variance analysis~(Section~\ref{sec:data:freq_analysis})
tells us exactly which bands to mix and which to leave alone.
Bins 0--8 ($\approx$36--325~Hz) are dominated by recording-condition
noise.
Bins 9--35 ($\approx$361--1{,}360~Hz) carry the mosquito wingbeat and
are species-discriminative.
So we mix \emph{only} bins 0--8 and protect the rest.

\textbf{How it works:}
FBS-Mix is an input-level batch transform, applied before the model
with probability $p{=}0.5$ per batch (stochastic application
prevents the model from depending on the augmentation).
It implements AdaIN~\cite{huang2017adain}---adaptive instance
normalisation---on the domain-noise band only.
Given a padded batch $X \in \mathbb{R}^{B \times T \times F}$, extract
the low-frequency band $X_{\mathrm{lo}} = X_{:,:,0:K}$ ($K{=}9$).
Sample a mixing coefficient $\lambda \sim \mathrm{Beta}(\alpha,\alpha)$
with $\alpha{=}0.1$ (a U-shaped distribution: usually $\lambda$ is near
0 or 1, so most batches see a clear style donor, and blending is
occasionally subtle) and a random batch permutation~$\pi$.
Compute per-sample, per-bin statistics $(\mu_i,\sigma_i)$ over
\emph{valid frames only}---padding frames are excluded, because a
63-frame clip padded to 200 frames is 68\% zeros and would severely
bias na\"ive global statistics.
Transfer the blended statistics to the band:
\begin{equation}
  \tilde{X}_{\mathrm{lo},i} =
    \underbrace{\bigl(\lambda\,\sigma_i + (1{-}\lambda)\,\sigma_{\pi(i)}\bigr)}_{\text{blended std}}
    \cdot \frac{X_{\mathrm{lo},i} - \mu_i}{\sigma_i}
    +\; \underbrace{\bigl(\lambda\,\mu_i + (1{-}\lambda)\,\mu_{\pi(i)}\bigr)}_{\text{blended mean}}.
  \label{eq:fbsmix}
\end{equation}
Bins $K{:}K'$ (wingbeat, $K'{=}36$, 361--1{,}360~Hz) and bins
$K'{:}F$ (upper harmonics and noise, 1{,}412--3{,}854~Hz) are left
unchanged.
Because FBS-Mix operates on the raw spectrogram before any model
layers, there is no inference-time overhead: it is simply switched
off during evaluation.

% =================================================================
\section{Experimental Setup}
\label{sec:experiments}
% =================================================================

\paragraph{Evaluation protocol.}
We use Leave-One-Domain-Out~(LODO): train on four of the five
recording conditions, evaluate on the fifth.
We run folds D1--D4 and report the mean; D5 is excluded because
holding it out leaves only 1,308 training clips---a data-starvation
problem~(Section~\ref{sec:data}).
All ablations use seed~42.
We separately run seeds 1234 and 3407 on the baseline and proposed
system to estimate variance; these are in Table~\ref{tab:multiseed}.

\begin{table}[h]
  \centering
  \caption{LODO D1--D4 mean $\BAu$ across seeds for baseline and
           proposed system. $^\star$D3 fold incomplete; 3-fold mean.}
  \label{tab:multiseed}
  \setlength{\tabcolsep}{5pt}
  \resizebox{\columnwidth}{!}{%
  \begin{tabular}{lccc}
    \toprule
    System & seed 42 & seed 1234 & seed 3407 \\
    \midrule
    MTRCNN baseline  & 0.166 & 0.197 & 0.147 \\
    Proposed (Bal+DANN+DiCL+proj128+$\tau$0.2) & 0.378 & $0.351^\star$ & 0.355 \\
    \bottomrule
  \end{tabular}%
  }
\end{table}

\paragraph{Metrics.}
\emph{Balanced accuracy}~(BA) is the mean per-class recall over the 9
species---it gives equal weight to rare species instead of letting the
numerically dominant species dominate the score.
Random chance for 9 classes is 11.1\%.
Following the challenge definition:
\begin{itemize}
  \item $\BAu$~(primary): BA on the held-out condition.
  \item $\BAs$: BA on the four seen conditions.
  \item $\DSG = |\BAs - \BAu|$: how much worse the model is on new
    conditions vs.\ known ones (lower is better).
\end{itemize}
All reported values are LODO D1--D4 mean.

\paragraph{Training.}
Adam optimiser, learning rate 1e-3, weight decay 1e-4, batch size 64.
We use early stopping with a minimum of 10 epochs, patience of~5
epochs, and a maximum of 100; the checkpoint with highest $\BAu$ on
the LODO validation split is used for evaluation.
Features are normalised by the per-bin mean and standard deviation
of the training split.
For DANN: $\lambda_{\max}{=}1.0$, $E{=}100$ (Eq.~\ref{eq:dann_schedule}).
For DiCL: weight $w_c{=}0.1$, temperature $\tau{=}0.2$, projection
head $32{\to}128{\to}128$ with ReLU between layers.
For FBS-Mix: $\alpha{=}0.1$, probability $p{=}0.5$, cutoff $K{=}9$.
All runs use one NVIDIA A100 GPU; each LODO fold takes roughly 27~min
for MTRCNN and 100~min for AST.

\paragraph{What we compare.}
We build up the system one component at a time to show the
contribution of each piece:
\begin{itemize}
  \item \textbf{Baseline}: unbalanced MTRCNN, standard cross-entropy.
  \item \textbf{+MixStyle}~/~\textbf{+DANN}: apply each to the
    \emph{unbalanced} data, to show they are near-useless without
    balance.
  \item \textbf{Balanced}: balanced resampling only.
  \item \textbf{Balanced + Species-only}$^*$: balanced, domain head
    turned off---tests whether domain supervision helps at all.
  \item \textbf{Balanced + DiCL}$^\dagger$: contrastive loss without
    adversarial training---tests whether DiCL works alone.
  \item \textbf{Balanced + DANN}: adversarial training on balanced data.
  \item \textbf{Balanced + DANN + DiCL}: add contrastive loss (32-dim,
    $\tau{=}0.07$, no projection head).
  \item \textbf{+proj128} / \textbf{+$\tau$0.2}: ablate the projection
    head and temperature.
  \item \textbf{Balanced + DANN + FBS-Mix}: test FBS-Mix as a
    DANN alternative (no DiCL).
  \item \textbf{Best + Augmentations$^\P$}: full proposed system plus
    HPSS and clip normalisation; results pending.
\end{itemize}
We additionally report two negative results in-text: test-time
batch-norm adaptation~(TTBN) and an Audio Spectrogram
Transformer~(AST) backbone.

% =================================================================
\section{Results \& Discussion}
\label{sec:results}
% =================================================================

\subsection{Main Ablation}
\label{sec:results:ablation}

Table~\ref{tab:ablation} shows LODO D1--D4 mean results (seed~42).

\begin{table}[t]
  \centering
  \caption{LODO D1--D4 mean results (seed~42). $\Delta$ = percentage-point
  gain in $\BAu$ over the unbalanced MTRCNN baseline (random chance 11.1\%).
  $^*$domain head disabled; $^\dagger$DiCL without DANN adds zero gain
  over balanced-only (explained in text);
  $^\P$results pending, updated before camera-ready; best in \textbf{bold}.}
  \label{tab:ablation}
  \setlength{\tabcolsep}{4.5pt}
  \renewcommand{\arraystretch}{1.08}
  \resizebox{\columnwidth}{!}{%
  \begin{tabular}{lcccc}
    \toprule
    Method & $\BAu$ & $\Delta$ & $\BAs$ & $\DSG$ \\
    \midrule
    \multicolumn{5}{l}{\textit{Without balanced resampling (near-useless)}} \\
    MTRCNN baseline & 0.165 & -- & 0.578 & 0.412 \\
    ~~+~MixStyle~\cite{mixstyle2021} & 0.167 & $+$0.2 & 0.592 & 0.425 \\
    ~~+~DANN~\cite{dann2016} & 0.169 & $+$0.4 & 0.621 & 0.453 \\
    \midrule
    \multicolumn{5}{l}{\textit{Balanced resampling: the prerequisite step}} \\
    Balanced & 0.321 & $+$15.6 & 0.511 & 0.245 \\
    Balanced + Species-only$^*$ & 0.324 & $+$15.9 & 0.547 & 0.267 \\
    Balanced + DiCL$^\dagger$ & 0.321 & $+$15.6 & 0.547 & 0.262 \\
    \midrule
    \multicolumn{5}{l}{\textit{Adding adversarial training on top of balanced data}} \\
    Balanced + DANN & 0.368 & $+$20.3 & 0.540 & 0.276 \\
    Balanced + DANN + DiCL & 0.340 & $+$17.5 & 0.517 & 0.247 \\
    Balanced + DANN + DiCL + proj128 & 0.371 & $+$20.6 & 0.540 & 0.235 \\
    \midrule
    \multicolumn{5}{l}{\textit{Proposed system (two operating points)}} \\
    Balanced + DANN + FBS-Mix & 0.345 & $+$18.0 & 0.516 & 0.229 \\
    \best{Balanced + DANN + DiCL + proj128 + $\tau$0.2} &
      \best{0.378} & \best{$+$21.3} & 0.534 & \best{0.202} \\
    ~~+ augmentations$^\P$ & \textcolor{pending}{---} & \textcolor{pending}{---}
      & \textcolor{pending}{---} & \textcolor{pending}{---} \\
    \midrule
    \multicolumn{5}{l}{\textit{$\ddagger$ Different protocol---not a fair comparison.
    Official 10-seed standard split:}} \\
    \multicolumn{5}{l}{\textit{~~all 5 conditions in training, test set 85\% D5-like.}} \\
    MTRCNN$^\ddagger$~\cite{biodcase2026} & 0.175 & -- & 0.881 & 0.706 \\
    \bottomrule
  \end{tabular}%
  }
\end{table}

\paragraph{Balanced resampling is the single most important step.}
Before any domain-generalisation method is applied, standard DANN
and MixStyle on the raw unbalanced data gain only 0.4 and 0.2~pp
respectively---barely above the 11.1\% random chance level.
Balanced resampling alone jumps to $\BAu{=}0.321$, a gain of 15.6~pp.
The mechanism is simple: the model was previously learning from one
recording condition 99.4\% of the time; now it sees all five equally.
This is not a DG method---it is just fixing the data imbalance that
was making every DG method ineffective.

\paragraph{DiCL without DANN does nothing (+0 pp).}
Balanced + DiCL (no DANN) achieves 0.321---identical to balanced-only.
Without adversarial pressure, the embedding still encodes
recording-condition information alongside species information.
DiCL pulls same-species cross-condition embeddings together, but the
species classifier simultaneously pushes different-species pairs
apart; on an embedding not yet cleared of condition-specific
directions, these two forces cancel and the net change is zero.
DANN first removes condition-specific directions from the embedding;
DiCL then has a cleaner space in which closeness means species
similarity rather than condition similarity.

\paragraph{DANN on top of balanced data adds +4.7 pp---but widens
the seen/unseen gap.}
Balanced + DANN reaches $\BAu{=}0.368$.
However, $\DSG$ actually increases from 0.245 to 0.276: the model
improves on unseen conditions, but also gets better on seen conditions
($\BAs$: 0.511 → 0.540), so the gap widens slightly.
This is a known tension in adversarial training: forcing the model to
ignore conditions helps on new ones but slightly over-removes
condition-correlated features that happen to be useful on known ones.
FBS-Mix and the DiCL projection head both reduce DSG on top of DANN.

\paragraph{DiCL needs a projection head and the right temperature.}
Adding DiCL directly on top of DANN (32-dim embedding, $\tau{=}0.07$)
\emph{decreases} $\BAu$ from 0.368 to 0.340.
Here is why: with $\tau{=}0.07$, the contrastive loss concentrates
almost all of its gradient on the single hardest negative in each
batch.
With only 80--634 minority-condition clips per epoch, most batches
have very few cross-condition same-species pairs.
The loss is trying to learn from sparse signal with a very sharp
gradient---essentially adding noise on top of the already-working
DANN signal.
A two-layer projection head ($32{\to}128{\to}128$) solves this by
learning the contrastive geometry in a separate, higher-dimensional
space without distorting the classifier's embedding.
Raising the temperature to $\tau{=}0.2$ spreads the gradient across
more negatives, producing a more stable training signal from the
small minority batches.
Together these recover +0.8~pp and reduce $\DSG$ further: best
result 0.378, $\DSG{=}0.202$.

\paragraph{FBS-Mix: a lighter-weight operating point.}
Balanced + DANN + FBS-Mix reaches $\BAu{=}0.345$, which is 2.3~pp
\emph{below} balanced + DANN alone.
FBS-Mix does not improve $\BAu$ over DANN; its value is in
$\DSG$: 0.229 vs.\ 0.276 for DANN, a reduction of 4.7~pp.
This makes FBS-Mix attractive when reducing the seen/unseen gap is
the priority, or when the projection head and contrastive loss are
undesirable (e.g.\ limited compute or a very small batch).
The contrast with MixStyle ($\approx$0.26~$\BAu$) confirms that
selective band mixing is the key: indiscriminate style mixing hurts
the wingbeat signal and gains nothing.

\paragraph{The official evaluation numbers are not comparable.}
The bottom row shows the official 10-seed baseline~\cite{biodcase2026}
for reference.
Those numbers come from a completely different setup: all five
conditions present in training, evaluated on a test set that is 85\%
D5-like.
$\BAs{=}0.881$ looks much higher than our LODO $\BAs$ because the
standard split trivially includes D5 in training.
$\DSG{=}0.706$ is so large because it is dominated by the D5-vs-rest
gap, not genuine cross-condition generalisation.
We include this row only to contextualise the gap between our
reported numbers and the challenge leaderboard; all method
comparisons in this paper use LODO exclusively.

\subsection{Why Standard Methods Fail: A Closer Look at Two Cases}
\label{sec:results:negative}

\paragraph{TTBN: adapting to the test set makes things worse.}
Test-time batch normalisation adaptation~(TTBN)~\cite{tent2021} puts
the batch normalisation layers back into training mode at inference
time, so each test batch updates the running mean and variance on the
fly rather than using the frozen training statistics.
Applied to our balanced+DANN checkpoint, TTBN \emph{harms} $\BAu$
by 3~pp.
The reason is straightforward: the test set is 85\% D5-like, so
each inference batch is dominated by D5 clips.
TTBN adapts the BN statistics toward D5, which then hurts recognition
of the D1--D4 clips that are the actual challenge.

\paragraph{AST: gradient reversal does not work on transformers.}
An Audio Spectrogram Transformer~(AST, $8{\times}8$ patches,
$d{=}192$, 4 layers, 1.8M params) trained from scratch reaches
$\BAu{=}0.214$~(unbalanced) and $\BAu{=}0.319$~(balanced)---similar
to MTRCNN (221K params).
But DANN \emph{fails} for AST: adding the GRL \emph{reduces} mean
$\BAu$ from 0.319 to 0.295, and per-fold results are highly
variable (D1: 0.109, D2: 0.448, D4: 0.334).
High variance across folds is the signature of a non-functional
adversarial signal.

Table~\ref{tab:ast_vs_mtrcnn} shows why.
For MTRCNN+DANN, training domain accuracy drops and stabilises at
0.28--0.30---just above the 5-class random-chance level of 0.20 under
balanced training.
The GRL is working: the model has been forced to forget which
recording condition it is in.
For AST+DANN, training domain accuracy stays at \textbf{0.997
throughout all epochs}.
The domain head trivially re-reads condition identity from the
transformer's output at every step; the gradient reversal never
wins.

\begin{table}[h]
  \centering
  \caption{Per-fold $\BAu$ (seed~42) and training domain accuracy.
           $^\dagger$Training domain accuracy under balanced training;
           5-class uniform chance = 0.20.
           MTRCNN+DANN reaches near-chance; AST+DANN stays at 0.997.}
  \label{tab:ast_vs_mtrcnn}
  \setlength{\tabcolsep}{4pt}
  \resizebox{\columnwidth}{!}{%
  \begin{tabular}{lcccc|c|c}
    \toprule
    Model & D1 & D2 & D3 & D4 & mean & tr.\ dom.\ acc$^\dagger$ \\
    \midrule
    MTRCNN baseline & 0.051 & 0.246 & 0.265 & 0.100 & 0.166 & 0.85 \\
    AST base        & 0.123 & 0.380 & 0.227 & 0.125 & 0.214 & 0.85 \\
    \midrule
    MTRCNN balanced & 0.186 & 0.629 & 0.343 & 0.127 & 0.321 & 0.73 \\
    AST balanced    & 0.204 & 0.522 & 0.311 & 0.236 & 0.319 & 0.86 \\
    \midrule
    MTRCNN+DANN & 0.271 & 0.739 & 0.328 & 0.134 & 0.368 & \textbf{0.30} \\
    AST+DANN    & 0.109 & 0.448 & 0.289 & 0.334 & 0.295 & \textbf{0.997} \\
    \bottomrule
  \end{tabular}%
  }
\end{table}

Two factors explain this.
First, AST's 1.8M-parameter backbone generates species-classification
gradients that are far larger than the small domain head's reversed
gradient---the adversarial signal is simply drowned out.
Second, MTRCNN's batch normalisation layers aggregate
recording-condition statistics across the \emph{batch}, concentrating
condition information in one learnable location that the GRL can
target directly.
AST's layer normalisation operates per-sample, so there are no
batch-level condition statistics to attack: condition information
is spread across all attention weights simultaneously, and the GRL
cannot localise it.
We are currently testing whether a 5$\times$ stronger GRL
($\lambda_{\max}{=}5.0$) or input-level FBS-Mix can rescue AST.

\subsection{Embedding Analysis}
\label{sec:results:tsne}

We extract the 32-dim embeddings by hooking the classifier input
and running t-SNE~(Fig.~\ref{fig:tsne}, pending GPU computation on
the LODO D1 checkpoints).
Earlier analysis on the standard-split baseline shows a clean
failure mode: all D1--D4 clips collapse into a single corner of the
embedding space, completely separate from D5.
The model routes every out-of-distribution input to the same place,
producing structured wrong predictions rather than calibrated
uncertainty.
The proposed system is expected to show well-mixed domain clusters,
and Fig.~\ref{fig:tsne} will be populated for camera-ready.

\subsection{Per-Species Analysis}
\label{sec:results:species}

On the official standard split, the baseline generalises almost
entirely through one species: \emph{Cx.\ pipiens}
($\BAu{=}0.804$).
This species has an unusually consistent wingbeat frequency across
recording conditions.
\emph{Cx.\ quinquefasciatus} and \emph{An.\ gambiae} collapse to
near-zero ($<$0.002), suggesting their wingbeat characteristics change
enough with recording condition to confuse the model completely.
Under LODO the same pattern holds.
The proposed system lifts most species into the 0.2--0.5~$\BAu$
range, with the largest absolute gains on species that are
under-represented in D1--D4, consistent with balanced sampling
providing their first meaningful training signal from minority
conditions.

% =================================================================
\section{Conclusion}
\label{sec:conclusion}
% =================================================================

The main lesson of this work is simple to state but easy to miss in
practice: if your training data is dominated by one condition, your
domain-generalisation method is not working---it is just learning
that one condition.
Fixing the imbalance first is not optional.
In the CD-MSC dataset, where one recording environment contributes
99.4\% of training clips, resampling to balance domains gains +15.6~pp
in cross-condition accuracy before any DG method is applied.
On top of that foundation, domain-adversarial training and
cross-condition contrastive learning push the system to $\BAu{=}0.378$
(+21.3~pp total) with a domain gap of $\DSG{=}0.202$.
FBS-Mix---which mixes only the recording-noise frequency band and
protects the wingbeat signal---provides a lightweight alternative that
reduces the seen/unseen gap to 0.229 with no model changes.

We also find that gradient reversal does not work for transformer
backbones in this setting: AST's domain classifier maintains 99.7\%
accuracy throughout training, meaning the adversarial signal never
suppresses condition information.
This failure comes from two compounding causes: the much larger
backbone gradient overwhelms the adversarial signal, and
per-sample layer normalisation removes the batch-level condition
statistics that gradient reversal relies on.

\paragraph{Limitations.}
All ablations use a single seed; full multi-seed variance estimates
for the proposed system are pending.
Domain labels are opaque, so we cannot say what physical differences
(device type, location, recording protocol) drive the cross-domain
gap.

\paragraph{Future work.}
The most promising direction for AST is applying domain-adversarial
training at the patch-embedding level rather than the output-token
level, where condition information may be more concentrated, or
using pretrained AST weights where some domain invariance is already
encoded.
For the MTRCNN system, combining FBS-Mix with the DiCL projection
head~(currently pending as ``best + augmentations'' in
Table~\ref{tab:ablation}) should reduce the seen/unseen gap below
0.200 while maintaining high $\BAu$.

% ------------------------------------------------------------------

\bibliographystyle{IEEEtran}
\bibliography{refs}

\end{document}
